% ============================================================
% RESULTS (part C): compiler performance, case studies, validity
% ============================================================
\subsection{Compiler Performance and Worked Examples (E5)}

For valid programs from 46 to 12{,}016 lines, a linear fit of total compile time against size gives \MsPerKloc{} ms per
thousand lines ($R^2=\ScaleRSq$), so a 40-line corpus program compiles in a few milliseconds. The log-log slope at large
sizes is \ScaleSlope, a little worse than linear. The cause is the semantic phase: it takes about a quarter of the time up
to 800 lines but 53\% at 12{,}000 lines, because every \code{if} copies the flow state, which grows with the number of
top-level declarations. Sharing those sets between branches would fix it, and at the size of real clinical scripts it
does not matter.

Across 300 generated programs the optimiser removed on average \OptMean\% of the three-address instructions (maximum
\OptMax\%), and the original and optimised code produced identical output on ten random inputs every time
(\OptMismatch{} mismatches in \OptPrograms{} programs). Generated programs have little redundancy, so the gain is modest.
On a demo program with constant arithmetic, the optimiser folds the constants, shares a repeated product
\code{per\_kg * weight} and deletes a dead branch, going from 18 instructions to 11. It still leaves a jump to the next
label (there is no peephole pass) and does not propagate across labels; neither affects correctness.

The command-line tool prints the compiler's diagnostics followed by the classifier's root-cause estimate.
Listing~\ref{lst:phi} shows a program with four faults: PHI concatenated into output, a temperature threshold missing
its decimal point (caught by both the range and the plausibility check), and a dose added to a weight. The classifier
ranks \emph{phi} first, which matches the root-cause priority.

\begin{lstlisting}[language={},numbers=none,basicstyle=\ttfamily\scriptsize,caption={Output of \texttt{medcheck.py examples/phi\_leak.med}.},label={lst:phi}]
== examples/phi_leak.med: 100 tokens, 4 error(s), 0 warning(s)
  12:19: error DOM002: protected health information reaches print() without anonymize()
  13:10: error DOM001: comparison against 380, outside the declared range [30, 43] of 'temp'
  13:12: error DOM001: value 380 degC is outside the plausible range [20.0, 45.0] degC
  16:12: error SEM005: unit mismatch: [mg] + [kg]
== learned root-cause estimate: phi 0.80, range 0.11, lexical 0.06
\end{lstlisting}

The more interesting case is a program the compiler has nothing to say about. In the example file
\texttt{inverted\_rule.med} the fever rule reads \code{if (temp < 38.0 degC)}. It compiles without a single
diagnostic, and the forest still puts it in \emph{logic} with probability 0.77 (valid 0.20). The tool prints a note
that the front end found no error but the classifier suspects a clinical-logic fault. It is worded as a suspicion
because it is an estimate from a model trained on synthetic faults, not a proof.

\subsection{Threats to Validity}
\label{sec:limits}

\begin{itemize}
    \item \textbf{Synthetic, in-distribution test.} Training and test programs come from the same generator and the same
    operators. The near-perfect accuracy shows how well the classifier learned this generator's faults. It says little
    about hand-written MedScript and nothing about real clinical code.
    \item \textbf{Labels come from the injection.} We checked that each fault shows up in the compiler, not that a human
    would call it the root cause. The ``most severe'' baseline is the labelling rule applied to diagnostics, so its 100\%
    on semantic families is partly circular.
    \item \textbf{Hand-written vocabulary.} The pattern features depend on a list of alert words, and a program that names
    its alerts differently would get past them.
    \item \textbf{No tuning.} Hyper-parameters were fixed in advance and never tuned, so a better-tuned GAN might do better.
    The scarcity range in E3 was widened after a pilot run on a separate, smaller corpus.
\end{itemize}
